\section{Methodology}
\label{sec:methodology}

\subsection{Sharing-ADMM coordination}
Let $x_i \in \mathbb{R}^{T}$ be the heating schedule of home $i$, $h_i$ its
baseline, and $B \in \mathbb{R}^{T}$ the fixed background aggregate. The
coordinator flattens the total feeder load toward a constant target $c$ while
respecting each home's comfort. We solve the sharing problem

\begin{equation}
\min_{\{x_i\}}\;
\sum_{i=1}^{N} \frac{\lambda}{2}\,\lVert x_i - h_i \rVert_2^2
\;+\; \frac{\mu}{2}\,\Big\lVert \sum_{i=1}^{N} x_i + B - c\,\mathbf{1} \Big\rVert_2^2
\label{eq:sharing}
\end{equation}

subject to~\eqref{eq:agent} for every $i$, where $c$ is the constant level that
conserves total daily energy, $\lambda = 0.15$ weights comfort deviation,
$\mu = 1.0$ weights flattening, and $\mathbf{1}$ is the all-ones vector. The
first term keeps each home close to its comfortable baseline; the second
penalizes any departure of the total load from a flat profile.

Following the sharing form of ADMM \cite{boyd2011admm}, we introduce an
auxiliary aggregate variable $z$ and a scaled dual $u$ with penalty
$\rho_{\mathrm{ADMM}} = 1.0$, and iterate three steps. The
\textbf{$x$-update} solves each home's local problem in parallel: a comfort-plus-penalty
proximal step followed by projection of the result onto the box-and-energy-plane
feasible set of~\eqref{eq:agent}. The \textbf{$z$-update} is the closed-form
average that drives the shared aggregate toward the flat target. The
\textbf{$u$-update} is the standard scaled dual ascent on the sharing residual.
Algorithm~\ref{alg:admm} summarizes one full sweep.

\begin{algorithm}[t]
\caption{Sharing-ADMM heating coordination}
\label{alg:admm}
\KwIn{baselines $\{h_i\}$, background $B$, target $c$, weights $\lambda,\mu$, penalty $\rho_{\mathrm{ADMM}}$}
\KwOut{coordinated schedules $\{x_i\}$}
initialize $x_i \leftarrow h_i$, $z \leftarrow 0$, $u \leftarrow 0$\;
\Repeat{aggregate residual converged}{
  \ForEach{home $i$ (in parallel)}{
    proximal step on $\frac{\lambda}{2}\lVert x_i - h_i\rVert^2$ plus penalty toward $z-u$\;
    project $x_i$ onto the box-and-energy set of~\eqref{eq:agent}\;
  }
  $z \leftarrow$ closed-form average enforcing the flat aggregate target\;
  $u \leftarrow u + \big(\textstyle\sum_i x_i + B\big) - z$\;
}
\Return{$\{x_i\}$}\;
\end{algorithm}

The $x$-update projection is what makes the scheme physically meaningful: it
guarantees that every iterate respects comfort and conserves daily energy, so the
coordinated schedule is feasible for the homes by construction. The number of
iterations to convergence is recorded as a tractability metric.

\subsubsection{Capped-energy projection in the $x$-update}
We make the local step explicit because its structure is what keeps every iterate
feasible. For home~$i$ the unconstrained comfort-plus-penalty objective is a
separable strictly convex quadratic, so its minimizer is available in closed form.
Writing $a_i = z - u$ for the current penalty anchor, the prox optimum is the
per-step convex combination
\begin{equation}
\tilde{x}_i \;=\;
\frac{\lambda\,h_i + \rho_{\mathrm{ADMM}}\,a_i}{\lambda + \rho_{\mathrm{ADMM}}},
\label{eq:prox}
\end{equation}
which trades the comfort pull toward $h_i$ against the consensus pull toward
$a_i$. This unconstrained point need not respect~\eqref{eq:agent}, so we
Euclidean-project it onto the intersection of the comfort box
$[(1-\alpha)h_{i,t},(1+\alpha)h_{i,t}]$ and the daily-energy hyperplane
$\sum_t x_{i,t} = \sum_t h_{i,t}$:
\begin{equation}
x_i \;=\; \arg\min_{x}\;\tfrac12\lVert x - \tilde{x}_i\rVert_2^2
\quad\text{s.t.}\quad
x \in \text{box}_i,\;\; \mathbf{1}^{\top}x = \mathbf{1}^{\top}h_i .
\label{eq:proj}
\end{equation}
Introducing a single scalar dual $\nu_i$ for the energy constraint, the projection
has the water-filling form
$x_{i,t} = \mathrm{clip}_{[(1-\alpha)h_{i,t},\,(1+\alpha)h_{i,t}]}(\tilde{x}_{i,t} + \nu_i)$,
and the value of $\nu_i$ that restores $\mathbf{1}^{\top}x_i = \mathbf{1}^{\top}h_i$
is found by a 1-D bisection on the scalar dual. Because the clipped sum is
monotone non-decreasing in $\nu_i$, the bisection is exact and costs only a handful
of evaluations, so the whole $x$-update is closed-form up to a scalar root-find.
The subsequent $z$-update is fully separable per time step --- each component of
$z$ is the average of the corresponding aggregate-plus-dual component driven toward
the flat target $c$ --- and therefore also closed form, which is why a full ADMM
sweep is inexpensive even with $N=30$ agents over $T=96$ steps.

\subsection{Imputation of non-responsive agents}
When agent $i$ is non-responsive its baseline $h_i$ is unavailable to the
coordinator, so we replace it with a learned estimate $\hat{h}_i$. We train a
single LightGBM gradient-boosted-tree regressor \cite{ke2017lightgbm} that maps a
small feature vector --- ambient temperature, the sine and cosine of the hour of
day, and the agent's nameplate heating level --- to instantaneous heating power.
The model is trained on the responsive population and applied step-by-step to
each silent agent, yielding $\hat{x}_{i,t}$ that the coordinator treats as the
agent's contribution. Generalization to unseen homes is assessed by a 5-fold
leave-agents-out cross-validation, which holds out entire agents rather than
random samples and therefore measures the realistic case of imputing a home the
model never saw.

\subsection{Power-flow validation loop}
The novel stage closes the loop on physics. For each scenario (uncoordinated,
ideal, and each $\varrho$) we take the resulting aggregate, distribute it across the
32 IEEE-33 load buses by their native load shares, and scale so that the
uncoordinated peak reaches the calibrated feeder. The trunk ampacity is fixed so
the uncoordinated peak loads the worst line to 115\%. We then run an AC
Newton-Raphson power flow at every one of the 96 time steps (672 solves across
all scenarios, all converged) and record, per scenario, the worst-of-day
minimum bus voltage and the worst-of-day maximum line loading. A scenario is
flagged as violating if any line exceeds 100\% loading at any step. This turns
an abstract kilowatt signal into a binary feeder verdict and a pair of
continuous severity measures.

Three modeling choices deserve comment. \emph{Load allocation by native share}
keeps the spatial signature of the IEEE-33 case intact, so the worst line is the
one the benchmark already stresses (the trunk leaving the substation) rather than
an artifact of our placement. \emph{Ampacity calibration to 115\%} is what makes
the experiment decision-relevant: a feeder that is comfortably within limits in
the uncoordinated case has nothing to gain from demand response, so we set the
binding constraint exactly where coordination must act to clear it. \emph{Solving
every step} rather than only the aggregate-peak step matters because imputation
error can shift the worst instant: the uncoordinated worst is the evening peak,
but under heavy imputation the binding step can migrate, and an
aggregate-peak-only check would miss it. Reporting the worst-of-day envelope
across all 96 steps therefore gives a conservative, instant-agnostic verdict.
